\chapter{Selected libraries}
\label{chap:Chapter6}

This chapter presents the selected Boost libraries for reimplementation in Rust.
The selection is based on several criteria, including their relevance to automotive software architectures, implementation characteristics, 
complexity and the potential benefits that a Rust implementation may provide in terms of memory safety, reliability and performance.

The selected libraries are: \texttt{Boost.Align} and \texttt{Boost.Endian}. Both libraries provide relatively low-level functionality
that is relevant to systems programming and embedded software, while also exposing functionality that interacts closely with memory
representation and processor architecture. These characteristics make them suitable candidates for evaluating the feasibility of 
reproducing existing C++ functionality using idiomatic Rust.

For each selected library, this chapter provides an overview of its purpose and functionality, describes its main components and discusses
its relevance to automotive software development. The chapter also identifies the main challenges associated with reimplementing the library
in Rust, including differences between the C++ and Rust memory models, language abstractions, portability requirements and interoperability
with existing C++ codebases.

\section{Boost.Align}
\subsection{Overview}

The \texttt{Boost.Align} library provides functions, classes, templates, traits and macros, for controlling, inspecting and diagnosing 
memory alignment. Its primary purpose is to provide portable alignment-related functionality across different compilers, standard library
implementations and language standards.

Memory alignment refers to the requirement that an object is stored at a memory address satisfying a particular alignment constraint. 
For example, a type with an alignment requirement of 8 bytes must generally be placed at an address divisible by 8. 
Alignment requirements are determined by the target architecture and by the type being stored. 
Correct alignment is important because processors may impose restrictions on the addresses from which particular types can be accessed. 
Even on architectures that support unaligned accesses, aligned accesses can provide more predictable and efficient memory operations.

Alignment is particularly relevant to low-level and performance-sensitive software. 
Incorrectly aligned objects may result in inefficient memory accesses or, on architectures with stricter alignment requirements, hardware faults.
In C++, violating the alignment requirements of an object can also result in undefined behavior. 
Consequently, explicit control over memory alignment can be important when implementing data structures, 
custom allocators, memory pools, serialization mechanisms, and other low-level components.

Before C++ 11, the language provided limited standardized facilities for expressing and manipulating alignment requirements.
Boost.Align addressed this limitation by providing portable alignment utilities that could be used in C++03 code and 
across implementations with incomplete or inconsistent support for newer alignment-related facilities.

C++11 introduced standardized support for specifying alignment requirements through language and library facilities. 
In particular, the language introduced alignment specifications through \texttt{alignas}, 
while the standard library introduced facilities such as \texttt{std::align}. 
However, compiler and standard library support was not always complete or consistent at the time Boost.Align was developed. 
The library therefore also provided portable implementations and compatibility facilities for environments in which 
the corresponding standard functionality was unavailable or unreliable.

\subsection{Functionality}

Boost.Align provides functionality for both dynamic memory allocation and pointer alignment.
Its facilities can broadly be divided into allocation, pointer alignment, alignment inspection, and alignment diagnostics.

For dynamic memory allocation, the library provides:

\begin{itemize}
\item \texttt{aligned\_alloc(alignment, size)}, which allocates a block of memory satisfying a specified alignment requirement;
\item \texttt{aligned\_free(ptr)}, which releases memory previously allocated using the corresponding aligned allocation mechanism.
\end{itemize}

These functions are useful when the default allocation provided by the runtime or standard library does not guarantee 
the alignment required by an application or data structure.

The library also provides allocator abstractions, including:

\begin{itemize}
\item \texttt{aligned\_allocator<T>};
\item \texttt{aligned\_allocator_adaptor<Allocator>};
\item \texttt{aligned\_delete<T>}.
\end{itemize}

The allocator abstractions allow alignment requirements to be integrated with C++ allocation mechanisms and standard containers.
This is particularly useful for containers whose elements must satisfy an alignment requirement greater than 
the default alignment provided by the allocator.

Another important part of Boost.Align concerns pointer alignment.
The library provides an implementation of functionality equivalent to \texttt{std::align}, 
allowing a pointer within an available memory region to be adjusted to satisfy a requested alignment 
while preserving the required amount of available storage.

This functionality is relevant when implementing custom memory-management mechanisms.
For example, an application may reserve a larger memory region and subsequently divide it into several aligned regions for different objects.
Explicit pointer alignment allows such regions to be placed at addresses compatible with the alignment requirements of the objects they contain.

The library additionally provides facilities for querying and inspecting alignment requirements.
These capabilities allow software to determine the alignment associated with a type or address and 
to verify whether an address satisfies a particular alignment constraint. 
Such facilities are useful when implementing generic low-level components where the alignment requirements cannot be assumed statically.

\subsection{Relevance to automotive software}

Memory alignment is particularly relevant to automotive software because many automotive systems operate under strict constraints 
regarding performance, memory consumption, determinism, and hardware compatibility. Software components may interact directly with hardware, 
communication buffers, memory-mapped data, DMA regions, or specialized data structures whose layout and alignment must be carefully controlled.

For example, an automotive application may process large collections of sensor data or 
communication frames using structures with specific alignment requirements. 
Correct alignment can contribute to efficient memory accesses and may be necessary to satisfy 
the requirements of a particular processor or hardware interface.

Alignment can also become important when software is designed to run across multiple processor architectures. 
Automotive platforms may use different CPU architectures and compiler toolchains, each with potentially different alignment characteristics. 
A portable abstraction for alignment can therefore reduce the amount of platform-specific code required by a software component.

Furthermore, automotive software frequently contains legacy C++ codebases that were developed before the widespread availability 
of modern C++ standards. Such systems may continue to rely on Boost libraries to provide functionality 
that was historically unavailable or insufficiently portable in the C++ language and standard library.

Although many of the capabilities historically provided by Boost.Align are now available through modern C++ language and 
standard library facilities, Boost.Align remains relevant when maintaining or migrating legacy systems. 
Reimplementing these facilities in Rust therefore provides an opportunity 
to investigate how low-level memory-management functionality originally designed for C++ can be expressed 
using Rust's stronger type system and ownership model.

\subsection{Challenges of reimplementation in Rust}

Reimplementing Boost.Align in Rust presents several challenges. 
The most significant is the interaction between explicit memory alignment and Rust's safety guarantees.

Rust provides alignment information as part of its type system and guarantees that references 
to properly typed values satisfy the corresponding alignment requirements. 
However, operations involving raw pointers, manually allocated memory, and custom allocation strategies 
may require the use of \texttt{unsafe} code. 
Consequently, an implementation of Boost.Align-like functionality must carefully isolate unsafe operations and 
ensure that the safety invariants expected by the Rust type system are preserved.

Dynamic aligned allocation is another important consideration. 
Rust's allocation APIs allow an allocation to be described through a \texttt{Layout}, 
which contains both the size and alignment of the requested memory region. 
This provides a natural abstraction for implementing functionality similar to Boost.Align's aligned allocation mechanisms. 
Nevertheless, the implementation must correctly handle invalid layouts, allocation failures, deallocation, 
and the relationship between an allocation's size and alignment.

Pointer alignment presents an additional challenge because it requires reasoning about raw memory addresses and available storage. 
A Rust implementation must ensure that pointer arithmetic does not violate the language's safety requirements and 
that any resulting pointer is only converted into a reference when all required invariants are satisfied.

Allocator integration also differs significantly between C++ and Rust. 
Boost.Align integrates with the C++ allocator model and standard containers, 
whereas Rust uses the \texttt{GlobalAlloc} and allocator-related abstractions available in the language ecosystem. 
Therefore, a direct translation of the C++ interfaces would not necessarily result in an idiomatic Rust implementation. 
Instead, the reimplementation should preserve the observable semantics of the original library while adapting its interface to Rust's ownership, 
borrowing, and allocation abstractions.

Finally, portability must be considered. 
The original Boost library was designed to compensate for differences between compilers and standard library implementations. 
A Rust implementation should therefore be evaluated across the relevant target architectures and 
toolchains to determine whether its behavior is consistent and whether platform-specific assumptions have been introduced.

\subsection{Rationale for selection}

Boost.Align was selected because it represents a low-level library whose functionality is closely related to memory management and 
processor architecture. It therefore provides a useful case study for evaluating the migration of C++ systems-level functionality to Rust.

The library is also particularly suitable for examining the boundary between safe and unsafe Rust. 
While many high-level uses of aligned data can be expressed safely through Rust's type system, 
explicit allocation and pointer manipulation require lower-level primitives. 
This makes Boost.Align useful for assessing whether Rust can provide a safer abstraction around functionality 
that traditionally requires careful manual reasoning in C++.

In addition, Boost.Align provides a relatively compact and well-defined API while still containing several distinct concepts, 
including alignment calculation, pointer manipulation, dynamic allocation, and allocator integration. 
This provides sufficient implementation complexity to evaluate the migration methodology without making the selected library unnecessarily large.

Although alignment support has substantially improved in modern C++ and much of the functionality 
provided by Boost.Align is now represented by standard C++ facilities, 
the library remains relevant when dealing with older C++ codebases and portability requirements. 
Its reimplementation in Rust can therefore demonstrate how existing low-level C++ infrastructure can be reproduced 
while taking advantage of Rust's memory-safety guarantees and modern abstractions.

The original Boost.Align documentation and source code were used as the primary references for the analysis and subsequent reimplementation:

\begin{itemize}
\item \url{https://github.com/boostorg/align/blob/boost-1.91.0/doc/align.qbk}
\item \url{https://www.boost.org/library/latest/align/}
\item \url{https://www.boost.org/doc/libs/latest/doc/html/align.html}
\end{itemize}

\section{Boost.Endian}
\subsection{Overview}

The \texttt{Boost.Endian} library provides facilities for manipulating the byte order, or endianness, of integer and 
user-defined types. Endianness defines the order in which the individual bytes of a multi-byte value are represented in
the memory. The two predominant byte-ordering schemes are little-endian, where the least significant byte is stored 
first, and big endian, where the most significant byte is stored first.

Endianness is generally transparent to application level software when the data remains within a single system.
However, it becomes relevant when binary data is exchanged between systems, stored in files, transmitted over networks,
or interpreted by software running on different processor architectures. In such cases, the representation of a value
in memory may differ from the representation expected by the receiving system.

As example, considering a 16-bit value such as \texttt{0x0102} is represented as the byte sequence \texttt{0x01 0x02} 
on a big-endian system and as \texttt{0x02 0x01} on a little-endian system. If the byte representation 
is interpreted without considering its endianness, the received application may obtain a different numerical value, 
leading to incorrect behavior or data corruption. Consequently, explicitly defining and handling byte order 
is an important requirement for portable binary data formats and communication protocols.

Boost.Endian was developed to provide a portable and consistent abstraction for these operations. The library supports 
binary data portability independently of the native endianness of the host platform and provides a common interface
across different C++ environments. It also addresses a secondary use case involving data size optimization by 
providing integer representations with sizes that are not directly available as standard C++ integer types, 
such as 24-bit, 40-bit and 48-bit integers. These types can be useful when implementing compact binary data formats or communication protocols.

The library is header only and in its current version, requires C++11. It can make use of compiler provided byte swapping intrinsics 
when available, allowing implementations to take advantage of architecture specific instructions while retaining a portable interface.

\subsection{Functionality}

Boost.Endian provides three complementary approaches for working with endian dependent data: endian conversion functions,
endian buffer types and endian arithmetic types. Each approach is designed for a different balance between explicitness, 
ease of use, performance and control over when conversions take place.

\subsubsection{Endian conversion functions}

The conversion function approach uses the standard C++ integer types to store values and explicitly converts their representation 
when required. The library provides both mutating and non-mutating conversion operations, including unconditional and conditional variants.

For example, an application can convert a value from big endian to the native byte order before performing calculations and convert
it back before writing the value to an external representation. This approach makes the conversion points explicit and is 
particularly appropriate when an application already has an established data model based on the native C++ types.

The conversion functions also provide a relatively direct abstraction over underlying processor byte swapping operations. When compiler
intrinsics are available, Boost.Endian can use them to generate efficient machine code.

\subsubsection{Endian buffer types}

Endian buffer types store values in a representation with a specified byte order. Unlike ordinary native integer types,
the representation maintained by the buffer does not change depending on the host architecture.

The library provides buffer types for a range of integer sizes, including 8-bit, 16-bit, 24-bit, 32-bit, 40-bit, 48-bit, 56-bit
and 64-bit integers. Both aligned and unaligned representations are available, with unaligned representations being particularly
useful when dealing with tightly packed binary formats.

An important characteristic of buffer types is that the conversions are explicit. A value must be explicitly loaded from or stored
into the endian representation. This prevents hidden conversions and gives the programmer precise control over when byte order 
transformation occurs.

This property is useful for binary protocols and file formats where the in memory representation must remain stable regardless 
of the architecture on which the software is executed.

\subsubsection{Endian arithmetic types}

Endian arithmetic types combine a fixed byte-order representation with an interface similar to the built-in C++ arithmetic types. 
They support arithmetic operations while automatically performing the required conversions between the stored representation and 
the native representation.

This approach simplifies application code because explicit conversions do not have to be performed for every arithmetic operation. 
For example, an application can maintain a structure containing big-endian integer fields and operate on those fields 
using normal arithmetic expressions.

The convenience of implicit conversion comes with a potential performance trade-off. 
Repeated operations on endian arithmetic types may result in repeated conversions, particularly inside computationally intensive loops. 
Boost.Endian therefore also documents patterns for moving conversions outside performance-critical loops when necessary.

\subsection{Binary data portability}

One of the primary motivations for Boost.Endian is binary data portability. 
Binary formats frequently specify the exact byte ordering and size of their fields. Unlike textual representations, 
binary formats cannot rely on the receiving system to interpret values according to its native representation.

This is particularly important for network protocols, file formats, serialization systems, and inter-process communication mechanisms. 
A data structure containing several integer fields may have to be transmitted from a little-endian processor to another system 
whose native representation is different. 
Explicitly representing the required byte order ensures that the same logical value is reconstructed at the destination.

Boost.Endian therefore separates the representation of externally defined binary data from the native representation of the host machine. 
This allows the application to use a consistent representation for data exchanged with external systems 
while remaining independent of the processor's native byte order.

The library also supports unaligned integer representations. 
This is useful for binary formats where fields are packed without the padding that would normally be introduced by C++ object alignment. 
Boost.Endian specifically identifies unaligned 24-, 40-, 48-, and 56-bit representations as useful 
when the exact size of the external representation matters.

\subsection{Relevance to automotive software}

Endianness is particularly relevant to automotive software because modern automotive systems consist of heterogeneous hardware 
and software components that communicate using precisely defined binary data formats.

Automotive systems commonly exchange data between electronic control units, sensors, gateways, communication interfaces, 
and higher-level computing platforms. These components may use different processors, operating systems, compilers, and software stacks. 
Consequently, communication interfaces cannot generally rely on the native representation of a particular processor.

The problem becomes particularly relevant when software processes binary communication frames or serialized data structures. 
A protocol may define a field as a 16-bit or 32-bit integer with a specific byte ordering, 
independently of the architecture executing the software. 
The implementation must therefore interpret the field according to the protocol definition rather than according to the host architecture.

This characteristic makes endian handling relevant to automotive communication stacks, serialization libraries, diagnostic protocols, 
and other components operating close to hardware and communication interfaces.

The relevance of these concerns is also reflected in modern automotive software platforms. 
The Eclipse S-CORE project, for example, is an open-source initiative focused on a software-defined vehicle platform 
and provides an example of the increasing role of modern software architectures and cross-platform technologies in automotive systems.

For an automotive software architecture, using explicit endian-aware representations can therefore help establish a clear boundary 
between the internal representation of application data and the representation required by an external communication or storage format.

\subsection{Memory representation and portability}

Boost.Endian is also relevant to portability because it avoids assuming that the host architecture uses the same byte order 
as the external data source.

A naïve implementation may interpret a binary structure directly through a native C++ structure. 
Such an approach can introduce portability problems because the resulting representation can depend on the processor's endianness, 
alignment requirements, padding rules, and compiler behavior.

For example, a structure containing several native integer fields cannot necessarily be treated as a portable representation 
of a network or file format simply by writing its memory contents directly to a file or communication channel. 
The representation of the fields may differ between platforms, and additional padding may be introduced by the compiler.

Endian-aware types provide a mechanism for explicitly representing the byte order required by the external format. 
This allows the application to distinguish between the representation used for communication and the representation used 
internally for computation.

This distinction is particularly valuable in safety- and reliability-oriented systems because it makes assumptions about 
binary representation explicit rather than relying on implicit properties of the target architecture.

\subsection{Performance considerations}

Byte-order conversion is generally a low-level operation that can be implemented efficiently using processor instructions. 
Boost.Endian takes advantage of compiler intrinsics when available, including byte-swap operations provided by GCC, Clang, 
and Microsoft Visual C++.

The performance characteristics of the different Boost.Endian approaches depend on how frequently conversions occur. 
When conversion functions are used, the programmer can explicitly control when a value is converted. 
This can be advantageous when a value is used repeatedly because the conversion can be performed once before a computationally intensive 
operation and the result converted back only when necessary.

Endian arithmetic types provide greater convenience but may perform implicit conversions during arithmetic operations. 
Boost.Endian's documentation demonstrates that, particularly in repeated computations, 
explicitly controlling the location of conversions can avoid unnecessary byte swapping.

This distinction is relevant for automotive applications because software may simultaneously have strict requirements for 
portability and computational efficiency. Communication-oriented code may prioritize explicit representation and correctness, 
while performance-critical processing may require minimizing repeated conversions.

\subsection{Challenges of reimplementation in Rust}

Reimplementing Boost.Endian in Rust presents several challenges, although Rust provides several primitives that are well suited to the problem.

The first challenge is representing values with an explicitly defined byte order while maintaining Rust's type 
and memory-safety guarantees. 
Rust's standard library provides methods such as \texttt{from\_be}, \texttt{from\_le}, \texttt{to\_be}, and \texttt{to\_le}, 
as well as byte-oriented conversion facilities. These provide a strong foundation for implementing endian-aware abstractions, 
but they do not directly reproduce the complete Boost.Endian API.

A second challenge is the implementation of endian-aware buffer types. 
Such types need to preserve a specific byte representation independently of the host architecture. 
A Rust implementation can use byte arrays as the underlying representation, which has the advantage of 
making the representation explicit and avoiding direct dependence on native integer layout.

Unaligned representations introduce another challenge. Rust provides strong guarantees around references and alignment, 
meaning that interpreting an arbitrary byte sequence directly as an integer reference may require unsafe operations. 
A safe implementation should therefore prefer byte-based representations and explicit conversion operations where possible, 
isolating any required unsafe code behind a well-defined abstraction.

A further challenge is reproducing Boost.Endian's arithmetic types. Boost.Endian allows endian-aware values to 
participate directly in arithmetic expressions, while Rust's operator traits provide a different abstraction mechanism. 
Implementing equivalent behavior requires implementing traits such as \texttt{Add}, \texttt{Sub}, \texttt{Mul}, 
and their assignment variants where appropriate.

However, a direct translation of the C++ API would not necessarily produce an idiomatic Rust library. 
Rust's ownership model, trait system, explicit conversion mechanisms, and emphasis on compile-time safety provide opportunities 
to design a safer interface rather than reproducing the C++ interface exactly.

Another important consideration is floating-point support. Endianness of floating-point representations cannot always be treated as 
a simple integer byte swap because floating-point representation and byte ordering involve additional assumptions. 
Boost.Endian therefore places restrictions on some floating-point conversion operations. A Rust reimplementation must 
preserve these semantic limitations rather than assuming that arbitrary floating-point values can safely be transformed by reversing their bytes.

\subsection{Rationale for selection}

Boost.Endian was selected because it represents a particularly relevant example of low-level, 
cross-platform functionality that is commonly required when implementing systems software.

Unlike functionality that is primarily concerned with convenience abstractions, 
endian handling directly affects the representation of data exchanged between software components. 
This makes the library particularly relevant to automotive architectures, where multiple software 
and hardware components must exchange binary data according to well-defined protocols.

The library also provides an interesting migration target because it contains several distinct implementation strategies: 
conversion functions, buffer types, and arithmetic types. These approaches expose different trade-offs between explicitness, 
ease of use, memory representation, and performance. 
Reimplementing them in Rust provides an opportunity to evaluate how these design choices can be expressed using Rust's type system 
and trait-based abstractions.

The support for both aligned and unaligned representations is another reason for its selection. 
It introduces a direct connection to low-level memory representation and requires careful consideration of alignment, 
raw byte access, and safe abstractions. This complements the selection of \texttt{Boost.Align}, which focuses explicitly on memory alignment.

The combination of \texttt{Boost.Align} and \texttt{Boost.Endian} is therefore particularly useful for the migration study. 
While Boost.Align focuses on the physical placement and alignment of objects in memory, 
Boost.Endian focuses on the representation and interpretation of multi-byte values. 
Together, they cover two closely related aspects of low-level systems programming: how data is positioned in memory and 
how that data is represented across system boundaries.

Furthermore, Boost.Endian is sufficiently mature and well-defined to provide a meaningful basis for 
comparing the original C++ implementation with a Rust implementation. The library's documented performance considerations, 
multiple API approaches, and portability requirements provide several dimensions against which the Rust implementation can be evaluated.

The primary references used for the analysis are:

\begin{itemize}
\item \url{https://www.boost.org/library/latest/endian/}
\item \url{https://www.boost.org/doc/libs/latest/libs/endian/doc/html/endian.html}
\item \url{https://github.com/eclipse-score}
\end{itemize}
